\newif\ifglobal

%%%%%%%%%%%%%%%%%%%%%%%
%%% Compile to view %%%
%%%%%%%%%%%%%%%%%%%%%%%

%%%%%%%%%%%%%%%%%%%%%%%%%%%%%%%%%%%%%%%%%%%%%%%%%%%
%%% Readme:                                     %%%
%%% https://www.overleaf.com/read/bwdjvfjksvjy/ %%%
%%%%%%%%%%%%%%%%%%%%%%%%%%%%%%%%%%%%%%%%%%%%%%%%%%%

% >>> M?: marks a module setting number or important notice

% >>> M4: Set {./relative-path} to external files for this module <<<
\def \modulefiles {./34-PCNT}
% To add an external file, use:
% (Replace '---' with actual filename)
% '\input{\modulefiles/tables/---.tex}' for tables
% '\includegraphics{\modulefiles/figures/---}' for figures
% '\subfile{\modulefiles/---}' for register files


% >>> M1: This module is in global project? <<<
% 'YES' - uncomment; 'NO' - comment out
\globaltrue

\ifglobal
    \documentclass[main\_\_CN.tex]{subfiles}

\else
    % Font "FandolSong-Regular" used by ctex does not have GSUB table.
    % It causes warnings that do not affect compilation and can be ignored.
    \PassOptionsToPackage{quiet}{fontspec}

    \documentclass[openany, 10pt]{book}

    % >>> M2: In ./00-shared/preamble-trm-module, also find
    %         settings for visibility of todo notes, watermark <<<
    \usepackage{preamble-trm-module}


    % Enable Chinese typesetting
    \usepackage{ctex}

    % >>> M3: Temporary labels in Chinese
    %         you can add your own labels there <<<
    \myexternaldocument{./00-shared/temp-labels-cn}

\fi %\ifglobal


\setcounter{secnumdepth}{4}


% >>> M5: If this module needs any updates in
% './00-shared/preamble-trm-module.sty'
% (include additional package, etc.), contact your module PM
% or create the file './00-shared/changelog.tex' make the
% necessary updates yourself and document them in this file <<<


\begin{document}


% Sets language into which pre-defined bits of text are translated
\selectlanguage{chinese}

% Do not delete this block
\ifglobal
% Add the following front matter if
% this module is outside of global project
\else
    \InputIfFileExists{readme}{}{}
    {\let\clearpage\relax \listoftodos}
    {\let\clearpage\relax \listofdonetodos}
    \tableofcontents
    %\listoftables
    %\listoffigures
\fi


%%%%%%%%%%%%%%%%%%%%%%%%%
%%% TRM Module Begins %%%
%%%%%%%%%%%%%%%%%%%%%%%%%

\hypertarget{pcnt}{}
\chapter{脉冲计数控制器 (PCNT)}
\label{mod:pcnt}


脉冲计数控制器 (Pulse Count Controller, PCNT) 用于对输入脉冲计数，通过记录输入脉冲信号的上升沿或下降沿进行递增或递减计数。PCNT 有四个称为“单元”的独立脉冲计数控制器，这些单元拥有自己的寄存器。PCNT 模块仅有一个时钟，为 APB\_CLK。下文描述中 \regindex{n} 表示单元编号 0 \verb+~+ 3。

每个单元有两个通道（ch0 和 ch1），可以独立配置为递增或递减计数。两个通道功能相同，下文以通道 0 (ch0) 为例进行介绍。

如图 \ref{fig:pcnt-block-diagram} 所示，每个通道有两个输入信号：
\begin{enumerate}
    \item 一个脉冲输入信号（如  sig\_ch0\_u\regindex{n} 为单元 \regindex{n} ch0 的脉冲输入信号）
    \item 一个控制信号（如 ctrl\_ch0\_u\regindex{n} 为单元 \regindex{n} ch0 的控制信号）
\end{enumerate}

\begin{figure}[H]
\includegraphics[width=0.4\textwidth]{\modulefiles/figures/pcnt-block-diagram.eps}
\centering
\caption{PCNT 框图}
\label{fig:pcnt-block-diagram}
\end{figure}

\section{主要特性}
PCNT 有如下特性：
\begin{itemize}
    \item 四个脉冲计数控制器（单元），各自独立工作，计数范围是 1 \verb+~+ 65535
    \item 每个单元有两个独立的通道，共用一个脉冲计数控制器
    \item 所有通道均有输入脉冲信号（如 sig\_ch0\_u\regindex{n}）和相应的控制信号（如 ctrl\_ch0\_u\regindex{n}）
    \item 滤波器独立工作，过滤每个单元输入脉冲信号（sig\_ch0\_u\regindex{n} 和 sig\_ch1\_u\regindex{n}）控制信号（ctrl\_ch0\_u\regindex{n} 和 ctrl\_ch1\_u\regindex{n}）的毛刺
    \item 每个通道参数如下：
        \begin{enumerate}
            \item 选择在输入脉冲信号的上升沿或下降沿计数
            \item 在控制信号为高电平或低电平时可将计数模式配置为递增、递减或停止计数
        \end{enumerate}
    \item 最大脉冲频率：$\frac{f_{APB\_CLK}}{2}$
\end{itemize}

\section{功能描述}\label{sec:pcnt-func-descr}

\begin{figure}[H]
\includegraphics[width=1\textwidth]{\modulefiles/figures/pcnt_arch_updated_20221026}
\centering
\caption{PCNT 单元基本架构图}
\label{fig:pcnt-arch}
\end{figure}

图 \ref{fig:pcnt-arch} 为 PCNT 单元的基本架构图。 如上所述， ctrl\_ch0\_u\regindex{n} 为单元 \regindex{n} ch0 的控制信号，控制信号 ctrl\_ch0\_u\regindex{n} 为高电平或低电平时可配置不同的计数模式，在输入脉冲信号 sig\_ch0\_u\regindex{n} 的上升沿和下降沿计数。可选计数模式如下：

\begin{itemize}
\item 递增模式：通道检测到 sig\_ch0\_u\regindex{n} 的有效边沿（软件可配）时，计数器的值 pulse\_cnt 加 1。pulse\_cnt 的值达到 \hyperref[fielddesc:PCNTCNTHLIMUN]{PCNT\_CNT\_H\_LIM\_U\regindex{n}} 时被清零。如果在 pulse\_cnt 达到 \hyperref[fielddesc:PCNTCNTHLIMUN]{PCNT\_CNT\_H\_LIM\_U\regindex{n}} 前， 该通道的计数模式改变或 \hyperref[fielddesc:PCNTCNTPAUSEUN]{PCNT\_CNT\_PAUSE\_U\regindex{n}} 置 1，则 pulse\_cnt 停止计数，计数模式改变。

\item 递减模式：通道检测到 sig\_ch0\_u\regindex{n} 的有效边沿（软件可配）时，计数器的值 pulse\_cnt 减 1。pulse\_cnt 的值达到 \hyperref[fielddesc:PCNTCNTLLIMUN]{PCNT\_CNT\_L\_LIM\_U\regindex{n}} 时被清零。如果在 pulse\_cnt 达到 \hyperref[fielddesc:PCNTCNTHLIMUN]{PCNT\_CNT\_L\_LIM\_U\regindex{n}} 前， 该通道的计数模式改变或 \hyperref[fielddesc:PCNTCNTPAUSEUN]{PCNT\_CNT\_PAUSE\_U\regindex{n}} 置 1， 则 pulse\_cnt 停止计数，计数模式改变。

\item 停止计数：计数停止，计数器的值 pulse\_cnt 保持不变。
\end{itemize}

表 \ref{tab:pcnt-positive-low} 至表 \ref{tab:pcnt-negative-high} 说明了如何配置通道 0 的计数模式。

\begin{table}[H]
\centering
\caption{控制信号为低电平时输入脉冲信号上升沿的计数模式}
\label{tab:pcnt-positive-low}
\begin{tabular}{|l|l|l|}
\hline
\rowcolor{lightgray}
\hyperref[fielddesc:PCNTCH0POSMODEUN]{PCNT\_CH0\_POS\_MODE\_U\regindex{n}} & \hyperref[fielddesc:PCNTCH0LCTRLMODEUN]{PCNT\_CH0\_LCTRL\_MODE\_U\regindex{n}} & 计数模式 \\ \hline
\multirow{3}{*}{1}             & 0                                     & 递增模式    \\ \cline{2-3}
                               & 1                                     & 递减模式    \\ \cline{2-3}
                               & 其他                                & 停止计数      \\ \hline
\multirow{3}{*}{2}             & 0                                     & 递减模式    \\ \cline{2-3}
                               & 1                                     & 递增模式    \\ \cline{2-3}
                               & 其他                                & 停止计数      \\ \hline
其他                         & N/A                                   & 停止计数      \\ \hline
\end{tabular}
\end{table}

\begin{table}[H]
\centering
\caption{控制信号为高电平时输入脉冲信号上升沿的计数模式}
\label{tab:pcnt-positive-high}
\begin{tabular}{|l|l|l|}
\hline
\rowcolor{lightgray}
\hyperref[fielddesc:PCNTCH0POSMODEUN]{PCNT\_CH0\_POS\_MODE\_U\regindex{n}} & \hyperref[fielddesc:PCNTCH0HCTRLMODEUN]{PCNT\_CH0\_HCTRL\_MODE\_U\regindex{n}} & 计数模式 \\ \hline
\multirow{3}{*}{1}             & 0                                     & 递增模式    \\ \cline{2-3}
                               & 1                                     & 递减模式    \\ \cline{2-3}
                               & 其他                                & 停止计数      \\ \hline
\multirow{3}{*}{2}             & 0                                     & 递减模式    \\ \cline{2-3}
                               & 1                                     & 递增模式    \\ \cline{2-3}
                               & 其他                                & 停止计数      \\ \hline
其他                         & N/A                                   & 停止计数      \\ \hline
\end{tabular}
\end{table}

\begin{table}[H]
\centering
\caption{控制信号为低电平时输入脉冲信号下降沿的计数模式}
\label{tab:pcnt-negative-low}
\begin{tabular}{|l|l|l|}
\hline
\rowcolor{lightgray}
\hyperref[fielddesc:PCNTCH0NEGMODEUN]{PCNT\_CH0\_NEG\_MODE\_U\regindex{n}} & \hyperref[fielddesc:PCNTCH0LCTRLMODEUN]{PCNT\_CH0\_LCTRL\_MODE\_U\regindex{n}} & 计数模式 \\ \hline
\multirow{3}{*}{1}             & 0                                     & 递增模式    \\ \cline{2-3}
                               & 1                                     & 递减模式    \\ \cline{2-3}
                               & 其他                                & 停止计数      \\ \hline
\multirow{3}{*}{2}             & 0                                     & 递减模式    \\ \cline{2-3}
                               & 1                                     & 递增模式    \\ \cline{2-3}
                               & 其他                                & 停止计数      \\ \hline
其他                         & N/A                                   & 停止计数      \\ \hline
\end{tabular}
\end{table}

\begin{table}[H]
\centering
\caption{控制信号为高电平时输入脉冲信号下降沿的计数模式}
\label{tab:pcnt-negative-high}
\begin{tabular}{|l|l|l|}
\hline
\rowcolor{lightgray}
\hyperref[fielddesc:PCNTCH0NEGMODEUN]{PCNT\_CH0\_NEG\_MODE\_U\regindex{n}} & \hyperref[fielddesc:PCNTCH0HCTRLMODEUN]{PCNT\_CH0\_HCTRL\_MODE\_U\regindex{n}} & 计数模式 \\ \hline
\multirow{3}{*}{1}             & 0                                     & 递增模式    \\ \cline{2-3}
                               & 1                                     & 递减模式    \\ \cline{2-3}
                               & 其他                                & 停止计数      \\ \hline
\multirow{3}{*}{2}             & 0                                     & 递减模式    \\ \cline{2-3}
                               & 1                                     & 递增模式    \\ \cline{2-3}
                               & 其他                                & 停止计数      \\ \hline
其他                         & N/A                                   & 停止计数      \\ \hline
\end{tabular}
\end{table}

每个单元均有一个滤波器，用于该单元的所有控制信号和输入脉冲信号。置位 \hyperref[fielddesc:PCNTFILTERENUN]{PCNT\_FILTER\_EN\_U\regindex{n}} 位使能滤波器。滤波器监测信号，滤除脉冲宽度小于 \hyperref[fielddesc:PCNTFILTERTHRESUN]{PCNT\_FILTER\_THRES\_U\regindex{n}} 个 APB 时钟周期的线路毛刺。

如前文所述，每个单元有通道 0 和通道 1 两个通道，处理不同的输入脉冲信号，并通过各自的 inc\_dec 模块递增或递减计数值。之后，两个通道将计数值发送给加法器 模块，该模块有一个带符号位的 16 位宽寄存器。软件可以通过置位 \hyperref[fielddesc:PCNTCNTPAUSEUN]{PCNT\_CNT\_PAUSE\_U\regindex{n}} 暂停加法器，也可以通过置位 \hyperref[fielddesc:PCNTPULSECNTRSTUN]{PCNT\_PULSE\_CNT\_RST\_U\regindex{n}} 清零加法器。

PCNT 可以设置五个观察点，五个观察点共用一个中断，可以通过每个观察点各自的中断使能信号开启或屏蔽中断。

\begin{itemize}
    \item 最大计数值：当 pulse\_cnt 大于等于 \hyperref[fielddesc:PCNTCNTHLIMUN]{PCNT\_CNT\_H\_LIM\_U\regindex{n}} 时， 产生上限中断， 同时 \hyperref[fielddesc:PCNTCNTTHRHLIMLATUN]{PCNT\_CNT\_THR\_H\_LIM\_LAT\_U\regindex{n}} 为高。
    \item 最小计数值：当 pulse\_cnt 小于等于 \hyperref[fielddesc:PCNTCNTLLIMUN]{PCNT\_CNT\_L\_LIM\_U\regindex{n}} 时，产生下限中断， 同时 \hyperref[fielddesc:PCNTCNTTHRLLIMLATUN]{PCNT\_CNT\_THR\_L\_LIM\_LAT\_U\regindex{n}} 为高。
    \item 两个中间阈值：当 pulse\_cnt 等于 \hyperref[fielddesc:PCNTCNTTHRES0UN]{PCNT\_CNT\_THRES0\_U\regindex{n}} 或者 \hyperref[fielddesc:PCNTCNTTHRES1UN]{PCNT\_CNT\_THRES1\_U\regindex{n}} 时，产生中断， 同时 \hyperref[fielddesc:PCNTCNTTHRTHRES0LATUN]{PCNT\_CNT\_THR\_THRES0\_LAT\_U\regindex{n}} 或 \hyperref[fielddesc:PCNTCNTTHRTHRES1LATUN]{PCNT\_CNT\_THR\_THRES1\_LAT\_U\regindex{n}} 为高。
    \item 零：当 pulse\_cnt 等于 0 时，产生中断， 同时 \hyperref[fielddesc:PCNTCNTTHRZEROLATUN]{PCNT\_CNT\_THR\_ZERO\_LAT\_U\regindex{n}} 有效。
\end{itemize}

在 PCNT 模块工作过程中，软件如果重新配置 \hyperref[fielddesc:PCNTCNTHLIMUN]{PCNT\_CNT\_H\_LIM\_U\regindex{n}} 和/或 \hyperref[fielddesc:PCNTCNTLLIMUN]{PCNT\_CNT\_L\_LIM\_U\regindex{n}}，则新的配置值会在 pulse\_cnt 达到上述五个观察点的任意一个后再生效；软件如果重新配置 \hyperref[fielddesc:PCNTCNTTHRES0UN]{PCNT\_CNT\_THRES0\_U\regindex{n}} 和/或 \hyperref[fielddesc:PCNTCNTTHRES1UN]{PCNT\_CNT\_THRES1\_U\regindex{n}}，则新的配置值会立即生效。

\section{应用实例}

每个单元的通道 0 和通道 1 可配置为独立工作或一起工作。下文详细说明了通道 0 独自递增计数、通道 0 独自递减计数和两个通道一起递增计数的应用实例。本节中未详述的通道工作模式（如通道 1 独自递减或递增、双通道一增一减），可参考这三种模式。

\subsection{通道 0 独自递增计数} \label{sec:pcnt-func-descr-one-channel-inc}

\begin{figure}[ht]
\fbox{\includegraphics[width=0.8\textwidth]{\modulefiles/figures/pulse_inc}}
\centering
\caption{通道 0 递增计数图}
\label{fig:pcnt-inc}
\end{figure}

图 \ref{fig:pcnt-inc} 为通道 0 在 sig\_ch0\_u\regindex{n} 上升沿独立递增计数的示意图，此时通道 1 关闭（请参阅 \ref{sec:pcnt-func-descr} 一节查看如何关闭通道 1）。通道 0 的配置如下所示。

\begin{itemize}
    \item \hyperref[fielddesc:PCNTCH0LCTRLMODEUN]{PCNT\_CH0\_LCTRL\_MODE\_U\regindex{n}}=0：当 ctrl\_ch0\_u\regindex{n} 为低电平时，递增计数。
    \item \hyperref[fielddesc:PCNTCH0HCTRLMODEUN]{PCNT\_CH0\_HCTRL\_MODE\_U\regindex{n}}=2：当 ctrl\_ch0\_u\regindex{n} 为高电平时，停止计数。
    \item \hyperref[fielddesc:PCNTCH0POSMODEUN]{PCNT\_CH0\_POS\_MODE\_U\regindex{n}}=1：在 sig\_ch0\_u\regindex{n} 的上升沿递增计数。
    \item \hyperref[fielddesc:PCNTCH0NEGMODEUN]{PCNT\_CH0\_NEG\_MODE\_U\regindex{n}}=0：在 sig\_ch0\_u\regindex{n} 的下降沿不计数。
    \item \hyperref[fielddesc:PCNTCNTHLIMUN]{PCNT\_CNT\_H\_LIM\_U\regindex{n}}=5：pulse\_cnt 的值递增至 \hyperref[fielddesc:PCNTCNTHLIMUN]{PCNT\_CNT\_H\_LIM\_U\regindex{n}} 时被清零。
\end{itemize}

\subsection{通道 0 独自递减计数} \label{sec:pcnt-one-channel-dec}

\begin{figure}[ht]
\fbox{\includegraphics[width=0.8\textwidth]{\modulefiles/figures/pulse_dec}}
\centering
\caption{通道 0 递减计数图}
\label{fig:pcnt-dec}
\end{figure}

图 \ref{fig:pcnt-dec} 为通道 0 在 sig\_ch0\_u\regindex{n} 上升沿独立递减计数的示意图，此时通道 1 关闭。此时通道 0 的配置与图 \ref{fig:pcnt-inc} 相比有如下区别：

\begin{itemize}
    \item \hyperref[fielddesc:PCNTCH0POSMODEUN]{PCNT\_CH0\_POS\_MODE\_U\regindex{n}}=2：即在 sig\_ch0\_u\regindex{n} 的上升沿递减计数。
    \item \hyperref[fielddesc:PCNTCNTLLIMUN]{PCNT\_CNT\_L\_LIM\_U\regindex{n}}=-5：pulse\_cnt 的值递减到 \hyperref[fielddesc:PCNTCNTLLIMUN]{PCNT\_CNT\_L\_LIM\_U\regindex{n}} 时被清零。
\end{itemize}

\subsection{通道 0 和通道 1 同时递增计数}\label{sec:pcnt-two-channels-inc}

\begin{figure}[ht]
\fbox{\includegraphics[width=0.9\textwidth]{\modulefiles/figures/pulse_add_2}}
\centering
\caption{双通道递增计数图}
\label{fig:pcnt-two-inc}
\end{figure}

图 \ref{fig:pcnt-two-inc} 为通道 0 和通道 1 分别在 sig\_ch0\_u\regindex{n} 和 sig\_ch1\_u\regindex{n} 上升沿同时递增计数的示意图。如图 \ref{fig:pcnt-two-inc} 所示，控制信号 ctrl\_ch0\_u\regindex{n} 与 ctrl\_ch1\_u\regindex{n} 的波形一致，输入脉冲信号 sig\_ch0\_u\regindex{n} 和 sig\_ch1\_u\regindex{n} 波形一致。具体配置如下：

\begin{itemize}
    \item 通道 0：
    \begin{itemize}
        \item \hyperref[fielddesc:PCNTCH0LCTRLMODEUN]{PCNT\_CH0\_LCTRL\_MODE\_U\regindex{n}}=0：当 ctrl\_ch0\_u\regindex{n} 为低电平时，递增计数。
        \item \hyperref[fielddesc:PCNTCH0HCTRLMODEUN]{PCNT\_CH0\_HCTRL\_MODE\_U\regindex{n}}=2：当 ctrl\_ch0\_u\regindex{n} 为高电平时，停止计数。
        \item \hyperref[fielddesc:PCNTCH0POSMODEUN]{PCNT\_CH0\_POS\_MODE\_U\regindex{n}}=1：在 sig\_ch0\_u\regindex{n} 的上升沿递增计数。
        \item \hyperref[fielddesc:PCNTCH0NEGMODEUN]{PCNT\_CH0\_NEG\_MODE\_U\regindex{n}}=0：在 sig\_ch0\_u\regindex{n} 的下降沿不计数。

    \end{itemize}

    \item 通道 1：
    \begin{itemize}
        \item \hyperref[fielddesc:PCNTCH1LCTRLMODEUN]{PCNT\_CH1\_LCTRL\_MODE\_U\regindex{n}}=0：当 ctrl\_ch1\_u\regindex{n} 为低电平时，递增计数。

        \item \hyperref[fielddesc:PCNTCH1HCTRLMODEUN]{PCNT\_CH1\_HCTRL\_MODE\_U\regindex{n}}=2：当 ctrl\_ch1\_u\regindex{n} 为高电平时，停止计数。

        \item \hyperref[fielddesc:PCNTCH1POSMODEUN]{PCNT\_CH1\_POS\_MODE\_U\regindex{n}}=1：在 sig\_ch1\_u\regindex{n} 的上升沿递增计数。

        \item \hyperref[fielddesc:PCNTCH1NEGMODEUN]{PCNT\_CH1\_NEG\_MODE\_U\regindex{n}}=0：在 sig\_ch1\_u\regindex{n} 的下降沿不计数。

    \end{itemize}
    \item \hyperref[fielddesc:PCNTCNTHLIMUN]{PCNT\_CNT\_H\_LIM\_U\regindex{n}}=10： pulse\_cnt 递增至 \hyperref[fielddesc:PCNTCNTHLIMUN]{PCNT\_CNT\_H\_LIM\_U\regindex{n}} 时被清零。

\end{itemize}


% >>> If your module has no registers, delete sample text
%       for Sections Register Summary and Registers below <<<

%%%%%%%%%%%%%%%%%%%%%%%%
%%% Register Summary %%%
%%%%%%%%%%%%%%%%%%%%%%%%

% >>> Update the sample text below and insert
%     the info on register summary into the subfile <<<

\hypertarget{pcnt-reg-summ}{}
\section{寄存器列表}
\label{sec:pcnt-reg-summ}

本小节的所有地址均为相对于 \textcolor{red}{脉冲计数控制器} 基地址的地址偏移量（相对地址），具体基地址请见章节 \ref{mod:sysmem} \textit{\nameref{mod:sysmem}} 中的表 \ref{tab:sysmem-base-address}。

请查看章节 \hyperref[glossary-access-types]{\textit{寄存器的访问类型}}，了解“访问”列缩写的含义。

\subfile{\modulefiles/34-PCNT-reg-summ__CN}


%%%%%%%%%%%%%%%%%
%%% Registers %%%
%%%%%%%%%%%%%%%%%

% >>> Update the sample text below and insert
%      the info on registers into the subfile <<<

\clearpage
\section{寄存器}

本小节的所有地址均为相对于 \textcolor{red}{脉冲计数控制器} 基地址的地址偏移量（相对地址），具体基地址请见章节 \ref{mod:sysmem} \textit{\nameref{mod:sysmem}} 中的表 \ref{tab:sysmem-base-address}。

\subfile{\modulefiles/34-PCNT-reg__CN}


\end{document}
